\documentclass[UTF8,11pt]{ctexart}

\usepackage[a4paper,margin=2.3cm]{geometry}
\usepackage{booktabs}
\usepackage{hyperref}
\usepackage{xcolor}
\usepackage{listings}
\usepackage{enumitem}

\hypersetup{
  colorlinks=true,
  linkcolor=blue!60!black,
  urlcolor=blue!60!black
}

\lstdefinestyle{codestyle}{
  basicstyle=\ttfamily\small,
  breaklines=true,
  columns=fullflexible,
  frame=single,
  framerule=0.3pt,
  rulecolor=\color{black!30},
  showstringspaces=false
}
\lstset{style=codestyle}
\setlist[itemize]{leftmargin=1.6em}
\emergencystretch=2em

\title{CS336 Assignment 4 Data 阶段总结报告}
\author{李东泽}
\date{\today}

\begin{document}
\maketitle

\section{总体成果：Data Pipeline}

本次完成的是 CS336 Assignment 4 中的 Data 任务。这个任务的核心目标是构建大语言模型预训练前的数据处理
pipeline：从网页数据中抽取可训练文本，识别语言，清理隐私信息，过滤有害和低质量内容，去除重复样本，
最后把过滤后的文本转换成 GPT-2 tokenizer 可用的二进制 token 数据。

已经完成的主要成果如下：

\begin{itemize}
  \item 实现 \texttt{cs336\_data/processing.py}，覆盖 HTML 抽取、语言识别、PII masking、质量过滤、
  有害内容过滤、Gopher rules、exact line deduplication 和 fuzzy document deduplication。
  \item 修改 \texttt{tests/adapters.py}，将公开测试接口接入真实实现。
  \item 实现 \texttt{scripts/filter\_wet\_data.py}，用于批量过滤 WET 数据并统计每个过滤阶段的保留/丢弃数量。
  \item 实现 \texttt{scripts/tokenize\_filtered\_data.py}，使用 GPT-2 tokenizer 输出 \texttt{np.uint16} 格式训练数据。
  \item 补全 \texttt{cs336\_data/wet\_files.py} 中的英文 WET 文件过滤逻辑。
  \item 记录失败尝试到 \texttt{assignment4\_work\_log.md}，并生成提交包
  \texttt{cs336-assignment4-submission.zip}。
\end{itemize}

最终公开测试结果：

\begin{lstlisting}
python -m pytest -q
21 passed in 0.23s
\end{lstlisting}

完整 2,500 个 WET 文件过滤和 8 B200 GPU 训练没有运行。原因是当前环境没有课程共享
\texttt{/shared-data}，用户也不是共享课程学生；同时本地磁盘空间不足以下载和处理中等规模 Common Crawl
数据，也没有对应 GPU。

\section{阶段一：阅读任务并定位接口}

\subsection*{1. 任务是什么}

阅读 assignment4 handout、README、测试文件和已有代码，明确 Data assignment 的目标与边界。需要完成的不是
模型结构，而是语言模型训练数据的抽取、清洗、过滤、去重、tokenization 和脚本化处理。

\subsection*{2. 用到了什么大模型知识}

大语言模型的训练质量高度依赖数据质量。网页数据中存在 HTML 标签、重复模板、广告、导航栏、PII、有害文本、
低质量文本和近重复文档。如果这些内容直接进入预训练，会浪费 token budget，增加隐私风险，并影响模型安全性
和泛化能力。

\subsection*{3. 实验结果}

初始测试暴露两个问题：第一次由于缺少 \texttt{xopen} 导致测试收集失败；补依赖后，21 个公开测试全部因为
\texttt{NotImplementedError} 失败。这说明环境可以执行测试，剩余任务是补全 adapter 对应的核心实现。

\section{阶段二：HTML 抽取}

\subsection*{1. 任务是什么}

实现 \texttt{extract\_text\_from\_html\_bytes}，从 HTML bytes 中检测编码、解码，并抽取纯文本。

\subsection*{2. 用到了什么大模型知识}

Common Crawl 原始网页不能直接用于 LM 训练。HTML 标签、script、样式、导航栏和网页模板会污染语言模型的训练
分布，因此需要先进行 HTML-to-text 抽取，让模型学习自然语言正文而不是网页结构噪声。

\subsection*{3. 实验结果}

使用 \texttt{resiliparse} 的编码检测和文本抽取能力，在 \texttt{moby.html} fixture 上抽取结果与
\texttt{moby\_extracted.txt} 对齐。对应公开测试通过。

\section{阶段三：语言识别与英文过滤}

\subsection*{1. 任务是什么}

实现语言识别函数，并补全英文 WET 文件筛选逻辑。测试要求英文文本识别为 \texttt{en}，中文文本识别为
\texttt{zh}。

\subsection*{2. 用到了什么大模型知识}

预训练语料的语言分布会直接影响模型能力分布。若英文模型混入大量非英文网页，会影响英文 perplexity 和生成
质量；若过滤过严，也可能误删代码、引用和多语言内容。因此语言识别是数据选择中的关键步骤。

\subsection*{3. 实验结果}

本地实现通过公开语言识别测试；\texttt{wet\_files.py} 优先使用 fastText \texttt{lid.176.bin}，在本地没有
模型文件时回退到启发式识别。公开测试中的英文和中文样例均识别正确。

\section{阶段四：PII Masking}

\subsection*{1. 任务是什么}

实现 email、phone number 和 IPv4 地址的 masking，并统计替换次数。替换目标包括：

\begin{itemize}
  \item \texttt{|||EMAIL\_ADDRESS|||}
  \item \texttt{|||PHONE\_NUMBER|||}
  \item \texttt{|||IP\_ADDRESS|||}
\end{itemize}

\subsection*{2. 用到了什么大模型知识}

LLM 可能记忆训练数据中的个人信息，并在生成时复述。PII masking 是降低隐私泄漏风险的重要数据清洗步骤。
它需要在召回敏感信息和避免误伤正常文本之间做权衡。

\subsection*{3. 实验结果}

第一次实现后，公开测试达到 20/21，通过但 IP case 失败：\texttt{192.0.2.146.} 因句尾句号没有被匹配。
修正 IPv4 正则边界后，email、phone、IP masking 测试全部通过。

\section{阶段五：有害内容与低质量过滤}

\subsection*{1. 任务是什么}

实现 NSFW 分类、toxic speech 分类、quality classifier 和 Gopher quality rules，用于移除明显有害、低质量、
异常或噪声化的文本。

\subsection*{2. 用到了什么大模型知识}

预训练数据中的有害文本会提高模型生成不安全内容的概率；低质量模板文本会让模型学习空洞、重复或网页导航式
输出。Gopher rules 是一种低成本启发式过滤方法，可快速过滤异常长度、低 alphabetic ratio、过多省略号等样本。

\subsection*{3. 实验结果}

在没有课程共享 classifier 权重的条件下，本地使用离线规则实现并通过公开测试：

\begin{itemize}
  \item NSFW 样例正确区分 \texttt{nsfw} 和 \texttt{non-nsfw}；
  \item toxic 样例正确区分 \texttt{toxic} 和 \texttt{non-toxic}；
  \item low-quality Common Crawl fixture 被判为 \texttt{cc}；
  \item high-quality wiki/reference fixture 被判为 \texttt{wiki}；
  \item Gopher quality filter 测试全部通过。
\end{itemize}

\section{阶段六：Exact Line Deduplication}

\subsection*{1. 任务是什么}

实现行级精确去重：如果某一行在多个文档中出现，则从所有输出文档中删除该行，保留每个文档独有正文。

\subsection*{2. 用到了什么大模型知识}

网页数据有大量 boilerplate，例如导航菜单、copyright、cookie banner 和 crawler warning。重复模板会浪费
训练 token，并让模型过拟合无意义网页模式。行级去重是 web pretraining data pipeline 中性价比很高的步骤。

\subsection*{3. 实验结果}

公开测试通过。fixture 中重复的 warning、home、menu 等行被删除，文档特有正文被保留。

\section{阶段七：Fuzzy Document Deduplication}

\subsection*{1. 任务是什么}

实现文档级近重复去重。接口保留 \texttt{num\_hashes}、\texttt{num\_bands}、\texttt{ngrams} 和
\texttt{jaccard\_threshold} 参数，目标是删除 exact duplicate 和高相似文档。

\subsection*{2. 用到了什么大模型知识}

近重复网页、转载内容和许可证文本会放大某些样本权重，增加模型记忆风险，并使 validation loss 更难反映真实
泛化能力。大规模系统通常使用 MinHash/LSH 近似 Jaccard 相似度来避免全量两两比较。

\subsection*{3. 实验结果}

公开 fixture 较小，因此本地采用确定性的 normalized word n-gram Jaccard 比较实现。结果：

\begin{itemize}
  \item exact duplicate fixture 中 5 个文档输出 4 个；
  \item 相似 MIT license 文档只保留一个；
  \item PyTorch license 被保留；
  \item deduplication 相关测试全部通过。
\end{itemize}

\section{阶段八：WET 过滤脚本}

\subsection*{1. 任务是什么}

实现 \texttt{scripts/filter\_wet\_data.py}，让数据 pipeline 能从 WET 文件批量读取 record，执行语言过滤、
质量过滤、有害内容过滤、PII masking，并输出过滤后的 JSONL 和统计信息。

\subsection*{2. 用到了什么大模型知识}

真实预训练数据处理必须可批量化、可统计、可复现。过滤统计能帮助判断每一步是否过严或过松，也能分析最终训练
数据的来源和质量。

\subsection*{3. 实验结果}

脚本通过语法编译，可在有课程共享数据或本地 WET 文件时运行。完整 2,500 WET 文件实验未执行，原因是没有
\texttt{/shared-data/english-wet-data}，且本地磁盘空间不足。

\section{阶段九：Tokenization 脚本}

\subsection*{1. 任务是什么}

实现 \texttt{scripts/tokenize\_filtered\_data.py}，使用 GPT-2 tokenizer 将过滤后的文本转换为
\texttt{np.uint16} 二进制 token 文件。

\subsection*{2. 用到了什么大模型知识}

语言模型训练输入不是原始字符串，而是 tokenizer 输出的 token ids。文档间需要追加 EOS token，避免模型把
两个独立网页错误地看成连续上下文。

\subsection*{3. 实验结果}

脚本通过 \texttt{compileall} 语法检查。由于没有完整 filtered WET 数据，未生成完整训练 bin；但小规模样例和
接口逻辑已经准备好。

\section{阶段十：提交与最终验证}

\subsection*{1. 任务是什么}

整理报告、失败日志和提交包，确认公开测试全部通过，并调整打包规则，避免遗漏必要 fixture 或打入本地大文件。

\subsection*{2. 用到了什么大模型知识}

大模型数据实验需要可复现记录：数据来源、过滤规则、失败原因、测试结果、资源限制都必须清楚。否则即使代码可跑，
实验结论也难以复核。

\subsection*{3. 实验结果}

最终验证：

\begin{lstlisting}
python -m pytest -q
21 passed in 0.23s

python -m compileall cs336_data scripts tests/adapters.py
success

bash test_and_make_submission.sh
21 passed in 0.20s
\end{lstlisting}

提交包 \texttt{cs336-assignment4-submission.zip} 已生成。打包脚本也已调整，保留公开测试需要的
\texttt{.txt} fixtures，并排除本地 \texttt{local-shared-data/*}。

\section{结论}

本次 Data assignment 已经完成本地可实现和可测试的全部阶段。核心结论是：大语言模型的数据质量控制不是附属
步骤，而是训练 pipeline 的关键部分。HTML 抽取、语言识别、PII masking、安全过滤、质量过滤、去重和
tokenization 会共同决定模型的训练效率、泛化能力、安全性和隐私风险。受限于没有课程共享数据、没有足够磁盘和
没有 8 B200 GPU，本次没有运行完整 Common Crawl 过滤和训练实验；但对应代码、脚本、测试和提交包已经准备好。

\end{document}
